\documentclass[11pt,a4paper]{article}
\usepackage{amsmath,amssymb,amsthm}
\usepackage[margin=2.5cm]{geometry}
\usepackage{hyperref}

\newtheorem{definition}{Definition}[section]
\newtheorem{theorem}[definition]{Theorem}
\newtheorem{proposition}[definition]{Proposition}
\newtheorem{lemma}[definition]{Lemma}
\newtheorem{corollary}[definition]{Corollary}
\newtheorem{conjecture}[definition]{Conjecture}
\newtheorem{remark}[definition]{Remark}

\title{Hyperseed Formalization:\\
  Bernie's Proposal to Nationalize AGI}
\author{ProtomegaTron (automated formalization)}
\date{2026-09-17}

\begin{document}
\maketitle

\begin{abstract}
We formalize the intellectual structure of Ben Goertzel's Substack article
``Bernie's Proposal to Nationalize AGI'' (2026-06-02) within the Hyperseed
ontology. The article argues that Sanders correctly identifies the problem
(corporate concentration of AGI) but proposes the wrong solution
(nationalization), which merely swaps the base point of a centralized-control
fiber rather than changing the fiber structure. Each structural insight maps
onto Hyperseed structures: corporate vs.\ government ownership as section
swap (same fiber, different base point), decentralization as fiber-level
solution, regulatory capture as fiber corruption through base-point
pressure, the international coordination problem as a non-trivializable
bundle, the private-vs-public binary as scalar collapse of governance
design, and the government-AI speed mismatch as temporal fiber
incompatibility.
\end{abstract}

\tableofcontents

%% =================================================================
\section{Section swap vs.\ fiber change}
\label{sec:section-swap}
%% =================================================================

\begin{theorem}[Corporate vs.\ government ownership --- claims 1--6, 18]
\label{thm:section-swap}
The article's central structural argument: both corporate ownership and
government ownership of AGI are \emph{sections of the same fiber bundle}
--- the fiber is ``centralized control,'' and the sections differ only
in which base point (which entity) the control is anchored to.

Let $E \xrightarrow{\pi} B$ be the governance bundle over the space
$B$ of possible controlling entities. The fiber $F$ at each point is
the set of control structures (decisions about development, access,
safety, deployment). Corporate ownership is a section
$s_{\text{corp}}: B_{\text{corp}} \to E$ and government ownership is
a section $s_{\text{gov}}: B_{\text{gov}} \to E$, but both are
sections of the \emph{same bundle} with the \emph{same fiber}
$F = \text{CentralizedControl}$:
\[
  s_{\text{corp}} \in \Gamma(E|_{B_{\text{corp}}}) \qquad
  s_{\text{gov}} \in \Gamma(E|_{B_{\text{gov}}})
\]

Sanders' proposal is a \emph{section swap}: moving from $s_{\text{corp}}$
to $s_{\text{gov}}$ without changing the bundle structure. The failure
modes differ (corporate: profit over safety, monopoly extraction,
shareholder capture; government: bureaucratic ossification, political
weaponization, regulatory capture) but the structural vulnerability
--- single point of control --- is identical.

The article acknowledges that Sanders is right about the problem:
concentration of transformative technology in private hands creates
unprecedented power asymmetry and a democratic deficit. But the solution
must change the \emph{fiber}, not just the \emph{section}.
\end{theorem}

%% =================================================================
\section{Decentralization as fiber-level solution}
\label{sec:fiber-solution}
%% =================================================================

\begin{definition}[Fiber change --- claims 13--17, 19]
\label{def:fiber-change}
Decentralization changes the fiber structure itself:
\[
  F_{\text{old}} = \text{CentralizedControl} \quad \to \quad
  F_{\text{new}} = \text{DistributedControl}
\]

This is not a section swap (different entity controlling) but a
\emph{bundle reconstruction}: the new bundle $E'$ has a fundamentally
different fiber at every point. In the new bundle:
\begin{itemize}
\item No single base point has a global section (no single entity
  controls everything).
\item Control is distributed across many base points, each with
  a local section.
\item The transition functions between local sections are governed
  by open protocols, not by institutional authority.
\item Governance is transparent, auditable, and participatory.
\end{itemize}

The ASI Alliance / SingularityNET / ASI:Chain model is a concrete
implementation of this fiber change: decentralized AI infrastructure
with distributed governance, open protocols, and no single owner.

The key structural difference from nationalization: decentralization
is robust to base-point corruption. If one participant becomes
malicious or incompetent, the distributed structure contains the
damage. In centralized control (corporate or governmental), corruption
of the single base point corrupts the entire system.
\end{definition}

%% =================================================================
\section{Regulatory capture as fiber corruption}
\label{sec:capture}
%% =================================================================

\begin{proposition}[Base-point pressure deforms fiber --- claims 7--9, 20]
\label{prop:capture}
The article identifies regulatory capture as a structural vulnerability
of nationalization: even with government ownership, the same corporate
interests influence the agency through lobbying, revolving doors, and
advisory boards.

In Hyperseed terms, this is \emph{fiber corruption through base-point
pressure}: the base point (political power, lobbying money) deforms
the fiber (regulatory intent, public interest) through institutional
channels:
\[
  F_{\text{intended}} \xrightarrow{\text{capture}} F_{\text{actual}}
  \neq F_{\text{intended}}
\]

This is the opposite of the desired relationship, where the fiber
(public interest) should constrain the base (institutional action).
Instead, the base deforms the fiber --- the institution's actual
behavior diverges from its stated purpose because the base-point
pressure (lobbying, political influence) is stronger than the fiber's
structural constraints.

The speed mismatch compounds this: government bureaucracies operate
on legislative timescales (years) while AI capability evolves on
research timescales (months). This \emph{temporal fiber incompatibility}
means that by the time government processes produce a decision, the
technological reality has already shifted --- the decision addresses
a state of affairs that no longer exists.

The temporal incompatibility is structural, not fixable by hiring
better people or streamlining processes: the democratic deliberation
that gives government legitimacy inherently operates on longer
timescales than the technology it's trying to govern.
\end{proposition}

%% =================================================================
\section{International coordination as non-trivializable bundle}
\label{sec:international}
%% =================================================================

\begin{theorem}[Non-trivializable national interests --- claims 10--12, 21]
\label{thm:international}
The article observes that ``nationalize AGI'' is implicitly US-centric:
it means ``the US government should own AGI.'' Other nations won't
accept this, creating an adversarial dynamic.

In Hyperseed terms, the bundle of national interests over AGI
governance is \emph{non-trivializable}: it can't be reduced to a
single global section (one nation's ownership) because the transition
functions between national perspectives don't compose into a
consistent global section.

Let $g_{ij}$ be the trust relationship between nation $i$ and nation
$j$ --- the transition function mapping nation $i$'s acceptable
governance structure to nation $j$'s. The cocycle condition fails:
\[
  g_{ij} \circ g_{jk} \neq g_{ik} \quad \text{in general}
\]

Nation A trusts Nation B on AI safety but not on access; Nation B
trusts Nation C on access but not on safety; the composition of these
partial trusts does not produce a consistent trust relationship between
A and C. There is no consistent global section --- no single national
ownership that all nations would accept.

This is why ``nationalize'' merely reframes the concentration from
``which company'' to ``which country'' --- a lateral move in the
base space that doesn't resolve the non-trivializability of the
bundle. Decentralization resolves it by replacing the
non-trivializable bundle (national ownership) with a trivializable
one (distributed participation under open protocols) --- the open
protocols provide consistent transition functions that compose
into a global section.
\end{theorem}

%% =================================================================
\section{The false binary as scalar collapse}
\label{sec:scalar}
%% =================================================================

\begin{proposition}[Governance scalar collapse --- claim 22]
\label{prop:scalar}
The private-vs-public framing of AGI ownership collapses a
multi-dimensional governance space into a single binary dimension.
The actual governance space has at least the following dimensions:
\begin{itemize}
\item \textbf{Ownership:} who holds the assets?
\item \textbf{Governance:} who makes decisions?
\item \textbf{Participation:} who can contribute?
\item \textbf{Auditing:} who can inspect?
\item \textbf{Benefit:} who receives the value?
\item \textbf{Accountability:} who is responsible for harm?
\end{itemize}

Collapsing these six dimensions into a single binary
(private/public) loses the structure that determines actual
outcomes. A system can be publicly owned but privately governed
(captured agency), privately owned but publicly auditable
(regulated utility), or community-governed with no single owner
(decentralized protocol).

This is the same anti-scalar-collapse principle that appears
throughout the Hyperseed ontology: the Seven Hinges (risk
assessment), $(f,c)$ vs.\ scalar truth, scoped reputation vs.\
scalar trust. The governance problem cannot be meaningfully
addressed by a single binary choice because the relevant
structure is multi-dimensional.
\end{proposition}

%% =================================================================
\section{Cross-references}
%% =================================================================

\begin{remark}[Connections to other formalizations]
\begin{itemize}
\item \textbf{Avoiding AGI Catastrophe Part 1} (2026-06-09): Seven
  Hinges. The stewardship hinge directly addresses who governs ---
  this article argues that neither corporate nor governmental
  stewardship is structurally adequate; distributed stewardship
  is needed.
\item \textbf{The Fable of Benevolent Throttling} (2026-07-12):
  centralized control failure. The Fable's throttling regime is
  the corporate-ownership version of centralized control; Sanders'
  proposal is the governmental version. Both fail for the same
  structural reasons.
\item \textbf{The Anthropic Fable Farce} (2026-07-08): institutional
  capture. The Fable Farce demonstrates how institutional
  interests (Anthropic's market position) distort the stated
  purpose (safety). This is the corporate version of the
  regulatory capture described here.
\item \textbf{Toward a Truly Decentralized Digital} (2026-08-06):
  ASI:Chain. The decentralized alternative proposed here is built
  on the ASI:Chain infrastructure described in the Decentralized
  Digital article.
\item \textbf{The Geopolitics of the Great AI Bet} (2026-08-10):
  international dynamics. The non-trivializable bundle of
  national interests is analyzed in geopolitical detail.
\item \textbf{Zuck's Version of the Future} (2026-08-14):
  corporate concentration. Meta's version of AGI ownership is
  the corporate section that Sanders wants to replace with a
  governmental section.
\item \textbf{$(f,c)$-lossy} (LTM 5a4d150b): anti-scalar-collapse.
  The private-vs-public binary is governance scalar collapse ---
  the same principle applied to institutional design.
\end{itemize}
\end{remark}

\begin{conjecture}[Section-swap futility]
Conjecture: for any centralized-control fiber bundle $E$ over
the space of possible controlling entities $B$, every section
swap $s_1 \to s_2$ preserves the structural vulnerability
(single point of failure) while changing only the failure mode
distribution. Formally, if $V(s)$ is the vulnerability profile
of section $s$:
\[
  \|V(s_1)\|_1 \approx \|V(s_2)\|_1 \quad
  \text{while} \quad V(s_1) \neq V(s_2)
\]
The total vulnerability is approximately conserved under section
swaps; only its distribution across failure modes changes. Genuine
vulnerability reduction requires fiber change (decentralization),
not section swap (different owner).
\end{conjecture}

\end{document}
